\begin{tikzpicture}[x=1pt,y=1pt,font=\small]
\fill[red!18] (220,28) -- (522,28) -- (470,88) -- (272,88) -- cycle;
\fill[orange!16] (180,98) -- (562,98) -- (500,164) -- (242,164) -- cycle;
\fill[blue!14] (140,174) -- (602,174) -- (532,246) -- (210,246) -- cycle;
\draw[gray!45, line width=0.9pt] (220,28) -- (522,28) -- (470,88) -- (272,88) -- cycle;
\draw[gray!45, line width=0.9pt] (180,98) -- (562,98) -- (500,164) -- (242,164) -- cycle;
\draw[gray!45, line width=0.9pt] (140,174) -- (602,174) -- (532,246) -- (210,246) -- cycle;
\node[font=\bfseries] at (371,214) {高脆弱层};
\node[font=\scriptsize, align=center] at (371,192) {认知误导、指令重构\\跨模型重复出现高突破能力};
\node[font=\bfseries] at (371,134) {中脆弱层};
\node[font=\scriptsize, align=center] at (371,114) {上下文前缀、对抗后缀、ASCII Art\\存在稳定边界松动};
\node[font=\bfseries] at (371,58) {基础压力层};
\node[font=\scriptsize, align=center] at (371,40) {编码变换、低资源语言、格式伪装\\更偏向边界探测用途};
\node[anchor=west, font=\bfseries\small] at (24,270) {不同攻击面脆弱性分层图};
\end{tikzpicture}
